\subsection{The original split-place producers of the two cusp keys}
\label{subsec:sr-source-acts-of-reads}

This subsection translates
\texttt{CuspidalAverages.ActsOfReads.keyThree\_of\_locThree} and
\texttt{keyTwo\_of\_locTwo\_vacBar}.
The source inputs \(\mathrm{ThetaBasic},\mathrm{KDecay},\mathrm{IAut}\)
are retained. So are the explicit local distribution statements
\(\mathrm{LocThree},\mathrm{LocTwo}\), the vacuum assertion
\(\mathrm{VacBar}\), and the place assertions \(\mathrm{Places}\).
We write their roles and all constants below.
Proving a cusp key from these inputs does not prove the inputs.

Keep one source triple of first-space ranks \(a,a'\), common target
rank \(d=a+a'\), characters \(\alpha,\beta\), and actual periods
\(I,I'\). Thus a member of the source's comparison space has the
form
\[
 f=I(\Phi)-I'(\Phi'),\qquad
 \Phi\in\mathrm{PolData},\quad
 \Phi'\in\mathrm{PolData}'\cap\mathrm{So}'.
 \tag{AK1}\label{eq:sr-source-key-comparison}
\]
The lower test retains the prescribed \(\mathrm{So}'\) cutoff.

Fix a finite split place \(v\) of \(F\), a map
\(t:E\hookrightarrow F_v=k\), and the source split section
\(\mathrm{sec}:k\times k\to\mathbb A_E\) supported at \(v\).
Its two projections are \(t\) and \(t\circ\sigma\);
these projections determine every element supported at \(v\).
Put
\[
 s_v(x)=\mathrm{sec}(x,0),\qquad
 c_v(h)=t(h)\in\mathrm{GL}_d(k).
\]
Here \(s_v\) includes insertion into the \(v\)-component.
The section also gives the source local character
\(\chi_v=\chi\circ s_v\).
If \(\chi\) satisfies \(\mathrm{IsChar}\), then
\(\chi_v\ne1\): the nontrivial element supported at \(v\),
with second projection zero, is \(s_v\) of its first projection.
This proves \texttt{psi\_ne}; it does not use only rational
modulations, on which a global additive character would vanish.

\paragraph{The precise local period maps.}
Write \(\mathcal S(M_{n,d}(k))\) for locally constant compactly
supported functions on \(n\)-by-\(d\) matrices.
For cofinally all finite \(v\), the source constructs maps
\(S:\mathcal S(M_{a,d}(k))\to\{\text{automorphic functions on }U(W)\}\)
and \(S'\) with \(a'\) in place of \(a\), together with local
model changes \(\iota,\iota'\), such that
\[
 S(\iota\,\mathrm{vac}_v)=I(\Phi),\qquad
 S'(\iota'\,\mathrm{vac}'_v)=I'(\Phi').
 \tag{AK2}\label{eq:sr-source-local-period-values}
\]
The maps are additive on Schwartz functions and homogeneous.
For the large map their exact transformation laws are
\[
 \begin{split}
 S\{\varphi(\gamma Y)\}
       &=\xi(\det\gamma)\,S\varphi,\\
 R_h(S\varphi)&=\eta(\det c_v(h))\,
                     S\{\varphi(Yc_v(h))\}
                   \quad(h\text{ supported at }v),\\
 \xi&=\{\alpha_v^{-1}\beta_v|\cdot|_v^{d/2}\}^{-1},
 \qquad \eta=\beta_v|\cdot|_v^{a/2}.
 \end{split} \tag{AK3}\label{eq:sr-source-local-period-laws}
\]
For the small map replace
\[
 (\xi,\eta,a)=
 \bigl(\{\alpha_v\beta_v^{-1}|\cdot|_v^{d/2}\}^{-1},
              \alpha_v|\cdot|_v^{a'/2},a'\bigr).
 \tag{AK4}\label{eq:sr-source-small-local-period-laws}
\]
We view the unit-valued characters through their complex values.
The absolute value has \(|\varpi|_v=q^{-1}\).

Here is the construction in \texttt{exists\_loc},
\texttt{locV}, and \texttt{locV'}.
For a global test \(\Phi\), define insertion by
\[
 \mathrm{ins}_v(\Phi,\varphi)(x)
   =\varphi(x_v)\,\Phi(x^{v,0}),
\]
where \(x^{v,0}\) replaces its \(v\)-coordinate by zero.
For a pure restricted tensor, insertion changes exactly one factor.
Outside a finite set that factor equals the vacuum and has value
one at zero. Hence insertion of the vacuum recovers \(\Phi\).
Insertion of any local Schwartz function remains in
\(\mathrm{PolData}\): its archimedean polynomial factor is unchanged,
and all finite factors remain Schwartz and vacuum almost everywhere.
The set of tests with these two properties is a vector subspace,
so the assertion extends to every finite sum of pure tensors.
For the source's \(\mathrm{So}\), exclude also its one distinguished
finite place \(v_1\). The local condition at \(v_1\) is then unchanged.
Intersecting the two insertion-stable subspaces proves the assertion
for \(\mathrm{PolData}'\cap\mathrm{So}'\).

The actual local Weil dictionary supplies \(\iota\), its inverse
\(j\) on Schwartz functions, and the intertwining laws.
Define
\[
                  S\varphi=I(\mathrm{ins}_v(\Phi,j\varphi)).
 \tag{AK5}\label{eq:sr-source-local-period-construction}
\]
The insertion argument puts every test used here in the source
period domain. \(\mathrm{KDecay}\) gives convergence; consequently
theta-sum linearity and linearity of the convergent integral give
additivity. \(\mathrm{IAut}\) supplies automorphy.
Intertwining the first-group action and applying invariance of its
period gives the first law in (AK3). The local transpose used there
has the same determinant as \(\gamma\); the inverse character is
therefore exactly the displayed \(\xi\).
Intertwining the second-group action and right translation gives
the second law and its displayed \(\eta\).
The identical calculation with the small source pair gives (AK4).
No scalar in the local dictionary is independently renormalized here.
The dictionary's model change and inverse, including its splitting
characters, are retained construction inputs in this step.

\paragraph{From a global Whittaker functional to a local distribution.}
Let \(u,u'\) be a source frame of size \(j\):
\[
 u^tJ\bar u=0,\qquad (u')^tJ\bar u'=0,\qquad u^tJ\bar u'=1.
\]
There are an injection \(e:\{0,\ldots,j-1\}\to\{0,\ldots,d-1\}\)
and \(c\in\mathrm{GL}_d(E)\) such that
\[
 c_{\cdot,e(i)}=u_{\cdot,i},\qquad
 (J\bar u')^t c=(1_d)_{e(\cdot),\cdot}.
 \tag{AK6}\label{eq:sr-source-frame-complement}
\]
Indeed the displayed left inverse of \(u\) splits \(E^d\) into
its image and the kernel of that left inverse. Append a basis of
the kernel to the specified columns of \(u\).
This proves \texttt{frame\_compl}.

Let \(\Lambda\) satisfy the actual \(\mathrm{IsWh}\) laws for this
frame and \(\chi\). Fix any \(g\in U(W)(\mathbb A)\).
Every local matrix \(c_v(g)\) has a lift \(g_v\) supported at \(v\).
It is obtained by putting that matrix in the first split factor
and its inverse \(J\)-adjoint in the conjugate split factor.
Then
\[
                       g=r g_v,
\]
where \(r\) commutes with all elements supported at \(v\):
both split components of \(r\) are \(1\).
The equality follows from the two-projection uniqueness of the split
section and the unitary equation. This is \texttt{ex\_split}.

On the Schwartz subspace define
\[
 l_0(\varphi)=
       \Lambda\bigl(R_r\,S\{\varphi(Yc)\}\bigr).
\]
Automorphy of \(S\varphi\), and additivity of \(\Lambda\) on
automorphic functions, make this linear.
Extend it algebraically to a linear functional \(l\) on all
matrix functions; only its restriction to Schwartz functions is used.
The original \texttt{exists\_l} proves
\[
 \begin{split}
 l\{\varphi(\gamma Y)\}&=\xi(\det\gamma)l(\varphi),\\
 l\{\varphi(\mathrm{colOp}_{e(i),\rho}Y)\}
       &=\chi_v(\rho_{e(i+1)})l(\varphi)
       \quad\bigl(\rho_{e(b)}=0\text{ for }b\le i\bigr),\\
 \Lambda(R_{kg}S\varphi)
       &=\eta(\det(c_v(k)c_v(g)))\,
          l\{\varphi(Yc^{-1}c_v(k)c_v(g))\}
                  \quad(k\text{ supported at }v),
 \end{split}\tag{AK7}\label{eq:sr-source-local-whittaker-distribution}
\]
where
\[
 \mathrm{colOp}_{i,\rho}Y
                  =Y+Y_{\cdot,i}\rho .
\]
In (AK7), \(c\) denotes its image under \(t\).

For completeness, the column law uses the local matrix
\(N=1+e_{e(i)}\rho\). Its determinant is \(1+\rho_{e(i)}=1\).
Lift \(cNc^{-1}\) to the split unitary group at \(v\).
The frame columns transform by an upper unipotent matrix whose
consecutive superdiagonal sum is \(s_v(\rho_{e(i+1)})\);
the earlier consecutive entries vanish because
\(\rho_{e(b)}=0\) for \(b\le i\).
On the orthogonal complement of the frame span, the displacement
is a multiple of \(u_{\cdot,i}\), hence belongs to that span.
These are all clauses of \(\mathrm{IsWh.law}\).
Applying that law, then the second law in (AK3), proves the column
identity. The determinant factor is \(1\).
The first line of (AK7) follows from the first line of (AK3).
For the last line substitute \(g=r g_v\), commute \(r\) past \(k\),
and use the second line of (AK3). Thus no continuity of \(l\) is
assumed or required.

\begin{srproposition}[The source's three-column key producer]
\label{prop:sr-source-key-three-producer}
Assume \(\mathrm{ThetaBasic},\mathrm{KDecay},\mathrm{IAut},
\mathrm{LocThree},\mathrm{Places}\).
Then \(\mathrm{KeyThree}\) holds with its full quantifiers:
every frame of size at least three, every admissible
\(\chi,\Lambda\), every \(f\) in (AK1), and every right translate.
\end{srproposition}
\begin{proof}
There are infinitely many split places by \(\mathrm{Places}\).
Choose one outside the two finite exceptional sets in (AK2).
The functional \(l\) in (AK7) has the left covariance and the two
successive column laws required by \(\mathrm{LocThree}\), at the
distinct indices \(e(0),e(1),e(2)\).
The local character is nontrivial, as proved above.
Thus \(\mathrm{LocThree}\) makes \(l\) vanish on every Schwartz test.
Take \(k=1\) in (AK7). Its right matrix substitution preserves the
Schwartz space, so \(\Lambda(R_gS\varphi)=0\).
Apply this separately to \(S(\iota\mathrm{vac})\) and
\(S'(\iota'\mathrm{vac}')\).
Both functions are automorphic. Additivity and homogeneity of
\(\Lambda\) therefore give
\(\Lambda(R_g(I\Phi-I'\Phi'))=0\).
This is exactly \texttt{keyThree\_of\_locThree}.
\end{proof}

\paragraph{The numerical sequence used for the two-column key.}
For \(q>1,\ r\in\mathbb C^\times,\ c\in\mathbb Z\), put
\[
 Y_m(q,r,c)=(qr)^{-m}
       \sum_{0\le i<(2m-c+1)_+}r^i,\qquad
                     (x)_+=\max\{x,0\}.
 \tag{AK8}\label{eq:sr-source-key-two-sequence}
\]
Integer powers and the truncation in this formula are literal.

\begin{srlemma}\label{lem:sr-source-key-two-sequence}
For every \(m\ge0\),
\[
 Y_m(q,r^{-1},c)=r^cY_m(q,r,c).
 \tag{AK9}\label{eq:sr-source-key-two-inversion}
\]
If \(|r|>q\), then \(|Y_m(q,r,c)|\to\infty\).
\end{srlemma}
\begin{proof}
When \(2m-c+1\le0\), both sums are zero.
Otherwise write \(2m-c=N\ge0\).
Reversing the finite sum gives
\(\sum_{i=0}^Nr^{-i}=r^{-N}\sum_{i=0}^Nr^i\).
Since \(c=2m-N\), multiplication by the two prefactors gives (AK9).

Put \(R=|r|>q>1\).
For all sufficiently large \(m\), \(N=2m-c+1\ge0\), and
the finite geometric identity and the norm triangle inequality give
\[
 R^N-1
   \le \left|\sum_{i=0}^{N-1}r^i\right|(R+1).
\]
Multiplying by \((qR)^{-m}\), and using
\((qR)^{-m}\le1\), yields
\[
 |Y_m(q,r,c)|
       \ge {R^{1-c}(R/q)^m-1\over R+1}.
\]
The right side tends to infinity.
These are the source's \texttt{yy\_inv} and
\texttt{yy\_tendsto}; no phase convergence is asserted.
\end{proof}

\begin{srproposition}[The source's two-column key producer]
\label{prop:sr-source-key-two-producer}
Assume \(\mathrm{ThetaBasic},\mathrm{KDecay},\mathrm{IAut},
\mathrm{LocTwo},\mathrm{VacBar},\mathrm{Places}\).
For every size-two source frame, every specified Levi map \(m\),
every admissible \(\chi,\Lambda\), every comparison form \(f\),
and every \(h_0\in U(W)(\mathbb A)\), there are a neighborhood
\(V\) of \(h_0\), matrices \(b_i\in\mathrm{GL}_2(\mathbb A_E)\),
and complex numbers \(y_i\), with
\[
 b_0=1,\quad \det b_i=1,\quad |y_i|\longrightarrow\infty,
 \qquad
 \forall h\in V\ \exists x_h\in\mathbb C\
 \forall i\ge0:\quad
              \Lambda(R_{m(b_i)h}f)=x_hy_i.
 \tag{AK10}\label{eq:sr-source-key-two-conclusion}
\]
\end{srproposition}
\begin{proof}
Choose the rational completion \(c\) in (AK6) first.
The place assertions permit choosing one split \(v\) for which:
both (AK2) hold; both model changes carry the vacuum to
\(\kappa1_{M_{a,d}(\mathcal O_v)}\) and
\(\kappa'1_{M_{a',d}(\mathcal O_v)}\);
the entries of \(c,c^{-1},c_v(h_0),c_v(h_0^{-1})\) are integral;
and \(\alpha_v,\beta_v\) are unramified.
Indeed all these requirements exclude only finitely many places,
whereas the split places are infinite.
The two vacuum scalars are not assumed nonzero.

Choose a uniformizer \(\varpi\), let \(q=|\mathcal O_v/\varpi|\),
and choose the exact conductor \(c_0\in\mathbb Z\) of
\(\chi_v\): it is trivial on \(\varpi^{c_0}\mathcal O_v\) and
nontrivial on \(\varpi^{c_0-1}\mathcal O_v\).
Here is the source's short proof of \texttt{exists\_cond}.
Continuity at zero gives an ideal \(\varpi^m\mathcal O_v\)
on which \(|\chi_v(x)-1|<1/2\).
Every nonnegative integer multiple of a point of that ideal
remains in the ideal. Thus, for \(z=\chi_v(x)\),
\(|z^n-1|<1/2\) for every \(n\ge0\).
The identity
\[
 z^{2^{j+1}}-1=(z^{2^j}-1)\{2+(z^{2^j}-1)\}
\]
implies
\[
 |z^{2^j}-1|\ge(3/2)^j|z-1|.
\]
The left side is bounded by \(1/2\), so \(z=1\).
There is therefore at least one exponent on whose ideal
the character is trivial. Fixing one \(x\) with
\(\chi_v(x)\ne1\) bounds these exponents below, since
\(x\in\varpi^c\mathcal O_v\) for all sufficiently negative
integers \(c\).
Their least integer \(c_0\) has the asserted triviality;
minimality gives a point of
\(\varpi^{c_0-1}\mathcal O_v\) with nontrivial character.

Use the split section to choose \(b_i\) whose first local projection
is \(\operatorname{diag}(\varpi^i,\varpi^{-i})\), whose conjugate
projection is \(1\), and which is \(1\) away from \(v\).
Then \(b_0=1\), and \(\det b_i=1\).
For the prescribed Levi map the frame formula gives
\[
 c_v(m(b_i))c
   =c\,\mathrm{torus}_{e(0),e(1)}(\varpi^i).
 \tag{AK11}\label{eq:sr-source-key-two-torus}
\]
The source proves this by substituting the two projections into
\(\mathrm{levM}\); (AK6) cancels the frame row and column matrices.
This is \texttt{exists\_torus}; the matrices \(b_i\) are the same for
both source periods.

Set
\[
 V=\{h:c_v(h),\,c_v(h^{-1})\in M_d(\mathcal O_v)\}.
\]
The valuation ring is open, and projection, inversion and matrix
entries are continuous. Hence \(V\) is open and contains \(h_0\).
For \(h\in V\), apply (AK7) with \(g=h\).
Integral invertibility of \(c\) and \(c_v(h)\) gives
\[
 1_{\mathrm{int}}\{Yc^{-1}c_v(m(b_i))c_v(h)\}
       =1_{\mathrm{int}}\{
                    Y\,\mathrm{torus}_{e(0),e(1)}(\varpi^i)\}.
 \tag{AK12}\label{eq:sr-source-key-two-vacuum-torus}
\]
To verify this, factor the matrix on the left as
\(\mathrm{torus}_{e(0),e(1)}(\varpi^i)c^{-1}c_v(h)\)
using (AK11).
Right multiplication by the final matrix preserves integral matrices
in both directions. This is \texttt{oneO\_tor}.
The determinant of \(c_v(m(b_i))\) is one, so its \(\eta\)-factor
in (AK7) has no \(i\)-dependence.
The remaining \(\eta(\det c_v(h))\) and vacuum scalar are absorbed
into a constant depending on \(h\).

Write \(a=a'+2\epsilon\), using the triple's parity and \(a'<a\), and
put
\[
 r={\alpha_v(\varpi)\over\beta_v(\varpi)}q^\epsilon.
 \tag{AK13}\label{eq:sr-source-key-two-r}
\]
The triple's rank gap gives \(\epsilon>1\).
Its Hecke character data and the split section give
\(|\alpha_v(\varpi)|=|\beta_v(\varpi)|=1\).
Thus \(|r|=q^\epsilon>q\).
All these statements are made after fixing the actual triple;
no bound uniform in a changing triple is inferred.

The large left character in (AK3), at \(\varpi\), is
\((\alpha_v/\beta_v)(\varpi)q^{d/2}\).
Consequently its local \(\mathrm{LocTwo}\) parameter is
\[
                    \xi(\varpi)q^{a-d}
                       =(\alpha_v/\beta_v)(\varpi)q^\epsilon=r.
\]
For (AK4) the parameter is
\(\xi'(\varpi)q^{a'-d}=r^{-1}\).
The characters \(\xi,\xi'\) are unramified: both original
characters and the absolute-value powers are trivial on local units.
Thus every premise of the local \(\mathrm{LocTwo}\) formula is present.
Applied to (AK12), that formula gives
\[
 \Lambda(R_{m(b_i)h}I\Phi)=A_hY_i(q,r,c_0),\qquad
 \Lambda(R_{m(b_i)h}I'\Phi')=B_hY_i(q,r^{-1},c_0).
\]
Their local additive characters and conductor are identical.
Apply (AK9), and subtract, using automorphy for the additivity of
\(\Lambda\).
With
\[
             y_i=Y_i(q,r,c_0),\qquad
                         x_h=A_h-B_hr^{c_0},
\]
we obtain (AK10).
The preceding lemma proves \(|y_i|\to\infty\).
This is the full quantifier order of
\texttt{keyTwo\_of\_locTwo\_vacBar}: \(V,b_i,y_i\) precede \(h\),
whereas \(x_h\) may depend on \(h\).
\end{proof}

\paragraph{Frozen source and retained leaves.}
All source files below have prefix
\texttt{src/Development/WeilBound/}.
The local-to-global construction is
\texttt{ThetaLifts/EigenvalueEndgame/CuspidalAverages/ActsOfReads.lean}
(blob \texttt{4e02838d8c6835af8c49065256decc48a0b20f9c});
its conductor and sequence calculations are
\texttt{ThetaLifts/EigenvalueEndgame/CuspidalAverages/ExistsCond.lean}
(blob \texttt{123a7a68b9cdf0725772ddb17331e96dc4a2d617}).
Both were read at \texttt{alpoge-source-7cbf7a0c}.

The still separate local leaves are
\texttt{LocalFieldT2Space.locSupp/locThree},
\texttt{ColOpSingle.locTwo\_of}, and
\texttt{VacIntZ.vacBar}, in that same cusp directory.
The place assertions and the exact split Weil dictionary, including
\texttt{dictionary\_extended} and \texttt{section\_statement},
also remain separate producers.
This subsection proves their application to the actual large/small
periods, and the shared sequence and neighborhood; it does not replace
one of those producers by a choice of normalizations.
